\subsection{Model description}

\paragraph{Migration as graph conductance.}
We represent a habitat by a connected, undirected graph $G=(V,E)$ with
$n=|V|$ nodes and $m=|E|$ edges. Let
$B\in\mathbb{R}^{m\times n}$ be an incidence matrix under an arbitrary edge
orientation. The target is a vector of relative edge log-conductances
$\vect{\theta}\in\mathbb{R}^m$. Positive conductances
$w_e=\exp(\theta_e)$ define the weighted graph Laplacian
\begin{equation}
    L(\vect{\theta})
    =B^\mathsf{T}\operatorname{diag}\{\exp(\vect{\theta})\}B.
    \label{eq:weighted-laplacian}
\end{equation}
The effective resistance between nodes $i$ and $j$ is
\begin{equation}
    R_{ij}(\vect{\theta})
    =
    (\vect e_i-\vect e_j)^\mathsf{T}
    L(\vect{\theta})^+
    (\vect e_i-\vect e_j),
    \label{eq:effective-resistance}
\end{equation}
where $L^+$ is the inverse on the mean-zero subspace. Multiplying every
conductance by the same constant changes the covariance scale but not the
relative surface. We therefore center true and estimated log-conductances to
have mean zero within each graph.

The canonical input is a node-level covariance matrix, coordinates for the
graph nodes, an undirected edge list, indices and sample counts for observed
nodes, and an effective SNP count. The public interface also accepts a sample
covariance matrix and sample coordinates, or genotypes and sample coordinates,
and constructs this object. Geographic coordinates assign samples to the
nearest node and determine normalized edge length. Environmental layers and a
fitted \feems surface are not inputs.

The model also predicts a centered node field
$\eta_i=\log q_i-\operatorname{mean}_j\log q_j$, where $q_i$ is the local
coalescence-rate coefficient used by the simulator. This nuisance output helps
the network distinguish variation in within-deme drift from edge conductance,
although the two remain statistically coupled.

\paragraph{Covariance calculation and deme aggregation.}
For a diploid genotype matrix $X$, with $s$ samples, $\ell$ SNPs, and entries
encoded as 0, 1, or 2, we remove monomorphic variants and standardize each
retained SNP across samples:
\begin{equation}
    Z_{al}
    =\frac{X_{al}-\bar X_{\cdot l}}
    {\operatorname{sd}(X_{\cdot l})},
    \qquad
    C=\ell^{-1}ZZ^\mathsf{T}.
    \label{eq:sample-covariance}
\end{equation}
The sample standard deviation uses denominator $s-1$. Missing genotypes in
the empirical analyses were mean-imputed after variant filtering. If $A$
averages samples assigned to each observed node, the node-level covariance is
$C_o=ACA^\mathsf{T}$.

When a user supplies covariance directly, we symmetrize it and project small
negative eigenvalues to zero before aggregation. The effective SNP count is
retained because it controls sampling-precision features. Covariance input
does not preserve linkage, allele-frequency spectra, or higher-order genotype
information.

\paragraph{Covariance score maps.}
We retain the leading 20 eigenpairs of the observed-node covariance and form
score maps
\begin{equation}
    U_k=\sqrt{\lambda_k}\vect v_k,
    \qquad k=1,\ldots,20.
    \label{eq:covariance-scores}
\end{equation}
Scores at unobserved nodes are zero, eigenvector signs are fixed
deterministically, and unavailable components are zero-padded. These are
eigenvectors of genetic covariance, not graph-Laplacian eigenvectors.

\subsection{All-pairs linearized resistance initializer}

The $o$ observed demes provide $o(o-1)/2$ pairwise dissimilarities
\begin{equation}
    D_{ij}=C_{o,ii}+C_{o,jj}-2C_{o,ij}.
\end{equation}
The initializer uses every observed pair. Its implementation batch size of
256 limits working memory and does not subsample the data.

We define a fixed baseline conductance
$w_e^{(0)}=\operatorname{clip}(d_{\mathrm{med}}/d_e,0.25,4)$ from edge length.
Sparse grounded solves give baseline resistances $R_{ij}^{(0)}$ and node
potentials $\vect\phi_{ij}$ for unit current injected at $i$ and removed at
$j$. The first-order sensitivity of log resistance to edge log-conductance is
\begin{equation}
    A_{(ij),e}
    =
    \left.
    \frac{\partial\log R_{ij}}{\partial\theta_e}
    \right|_{\vect\theta^{(0)}}
    =-
    \frac{w_e^{(0)}\{(B\vect\phi_{ij})_e\}^2}
    {R_{ij}^{(0)}}.
    \label{eq:resistance-sensitivity}
\end{equation}
Thus a pair contributes to each edge carrying baseline electrical flow between
its endpoints.

Let $\vect r$ be the centered difference between observed and baseline log
dissimilarity. We accumulate $\vect g=A^\mathsf{T}\vect r$ and the diagonal
information approximation $h_e=\sum_{i<j}A_{(ij),e}^2$ without materializing
the dense sensitivity matrix. The update solves
\begin{equation}
    \left[
      \operatorname{diag}(\vect h+\rho)+0.01L_E
    \right]\vect\delta=\vect g,
    \label{eq:linearized-initializer}
\end{equation}
where $L_E$ is the unweighted line-graph Laplacian and $\rho$ is 0.001 times
the median positive diagonal information. We center and scale $\vect\delta$
to unit standard deviation. The field and a numerical-validity indicator are
the final two edge features. This is a one-step diagonal approximation. It is
evidence for the network to denoise rather than a stand-alone inverse
solution.

\subsection{Fixed diffusion-band features}

The spatial representation is unchanged by translating, rotating, reflecting,
uniformly rescaling, or relabeling the same graph. The network receives no
absolute coordinates or signed $x$/$y$ edge displacement. Coordinate slots
retained for checkpoint-shape compatibility are fixed at zero. Edge length is
divided by a graph-level reference length, endpoint sample-count and precision
features use symmetric means and absolute differences, and support is routed
over all shortest paths rather than one arbitrarily selected path.

For the multiscale features, the reference edge weight is
$a_e=\operatorname{clip}(d_{\mathrm{med}}/d_e,0.25,4)$. If $A_0$ and $D_0$
are the corresponding adjacency and degree matrices, the fixed normalized
Laplacian is
\begin{equation}
    L_0=I-D_0^{-1/2}A_0D_0^{-1/2}.
\end{equation}
This operator depends on graph topology and relative edge length, not on the
unknown migration surface. For a node-signal matrix $S$, we compute
$S_t=\exp(-tL_0)S$ at $t=(0.25,0.5,1,2,4)$ and form
\begin{equation}
 \begin{split}
    \Psi_0(S)&=S-S_{0.25},\\
    \Psi_k(S)&=S_{t_k}-S_{t_{k+1}},\quad k=1,\ldots,4,\\
    \Psi_5(S)&=S_{4}.
 \end{split}
 \label{eq:diffusion-bands}
\end{equation}
The bands retain changes removed at successive diffusion scales and a final
coarse component.

The diffusion input contains 13 standardized node signals: the first eight
covariance score maps, log sample count, a sampled-node indicator, total and
within-deme covariance support, and the covariance diagonal. Six bands give
78 node features. For each signal and band, we compute three symmetric edge
summaries: absolute endpoint contrast, log one plus squared contrast, and mean
endpoint absolute energy. These give 234 edge features. Base features add the
full 20 covariance scores, sample and SNP precision, covariance support,
normalized length, symmetric endpoint summaries, and shortest-path support.
The final arrays have 109 node and 253 edge columns, including the two
initializer columns. Training-corpus means and standard deviations normalize
each column, with a standard-deviation floor of 0.2.

\subsection{Graph neural network}

Table~\ref{tab:model-setup} summarizes the model. We implemented it with
Equinox and JAX \citep{jax2018,deepmind2020jax,equinox2021}. An affine GELU
layer embeds node features in 128 dimensions. Ten residual edge-aware
message-passing layers apply directed messages
\begin{equation}
    \vect m_{j\rightarrow i}^{(k)}
    =
    \vect h_j^{(k)}\odot
    \operatorname{sigmoid}(W_g^{(k)}\vect x_{ij})
    +\operatorname{GELU}(W_m^{(k)}\vect x_{ij}),
    \label{eq:edge-message}
\end{equation}
followed by a residual node update using the mean incoming message. A
three-layer context branch uses the same graph. No graph coarsening or pooling
is present in the released architecture.

For edge $(i,j)$, the decoder uses
\begin{equation}
    \vect u_{ij}
    =
    [\vect h_i+\vect h_j,
      |\vect h_i-\vect h_j|,
      \vect x_{ij}],
\end{equation}
which is symmetric in its endpoints. A query derived from $\vect u_{ij}$
attends to keys and values from all context-node embeddings, giving every edge
graph-wide context despite finite message-passing depth. Separate affine heads
predict an edge mean $\mu_e$ and log variance $\ell_e$, with
$\ell_e\in[-7,5]$, and analogous node-$q$ outputs. Edge means and node means
are centered within each graph.

The output family is a factorized normal distribution trained against
simulated targets. It approximates the conditional distribution induced by
the simulation corpus. It is not the result of variational inference for an
explicit per-dataset likelihood, and it does not represent covariance among
edge errors. We therefore call $\exp(\ell_e/2)$ the network predictive SD, not
a posterior SD.

\begin{table}[t]
\centering
\caption{Model and optimization settings for the production checkpoint.}
\label{tab:model-setup}
\small
\begin{tabular}{ll}
\toprule
Component & Setting \\
\midrule
Node / edge dimensions & 109 / 253 \\
Hidden channels & 128 \\
Local / context graph layers & 10 / 3 \\
Graph-wide attention & yes \\
Covariance score maps & 20 \\
Diffusion times & 0.25, 0.5, 1, 2, 4 \\
Spatial representation & topology and normalized edge length \\
Coordinate-frame invariant & translation, rotation, reflection, scale \\
Linearized initializer & all\_pair\_diagonal\_adjoint \\
Initializer pairs & all observed-deme pairs \\
Predictive family & factorized\_normal \\
Training examples & 61,440 \\
Validation / held-out examples & 1,024 / 4,096 \\
Best / final step & 29,000 / 41,000 \\
Best validation loss & 0.1572 \\
Optimizer / learning rate / batch & Adam / 5e-5 / 8 \\
\bottomrule
\end{tabular}

\end{table}

\subsection{Graph-balanced training objective}

Let $\theta_e^\star$ and $\eta_i^\star$ be centered simulation targets. The
principal edge term is the graphwise heteroscedastic Gaussian negative
log-likelihood
\begin{equation}
    \mathcal L_{\mathrm{edge}}
    =\frac{1}{2m}\sum_{e=1}^m
    \left[(\mu_e-\theta_e^\star)^2\exp(-\ell_e)+\ell_e\right].
    \label{eq:edge-nll}
\end{equation}
The node-$q$ head uses the same form. The full loss is
\begin{equation}
 \begin{split}
    \mathcal L={}&
    \mathcal L_{\mathrm{edge}}
    +0.5\mathcal L_{\mathrm{context}}
    +0.5\mathcal L_q\\
    &+0.25\mathcal L_{\mathrm{slope}}
    +0.20\mathcal L_{\mathrm{corr}}
    +0.35\mathcal L_{\mathrm{orientation}}.
 \end{split}
 \label{eq:total-loss}
\end{equation}
The context term is edge MSE from the auxiliary branch. The slope term is a
Huber loss on the within-graph regression slope relative to one, the
correlation term is $1-r$, and the orientation term penalizes failure to
separate above- and below-mean target edges. Shape terms are disabled when the
target's centered RMS is at most 0.05. Every component is reduced within each
graph before averaging over the minibatch, so a graph contributes one weight
regardless of its edge count. The training loss contains no resistance or
covariance-forward term.

\subsection{Synthetic training corpus}

\paragraph{Graphs, surfaces, and sampling.}
We generated 61,440 training examples, 1,024 validation examples, and 4,096
held-out examples with disjoint random seeds. Table~\ref{tab:training-corpus}
reports the realized training ranges. The corpus combines five graph families.

The largest component contains 34,605 random global-landscape examples. We
sample geographic windows over land between $55^\circ$S and $72^\circ$N from
Natural Earth 50-m polygons \citep{naturalearth}, project each window to a
local azimuthal-equidistant system, and place a hexagonal mesh over its full
bounding box. Water nodes remain in the graph but cannot be sampled. Target
sizes are 220, 450, 900, or 1,400 nodes. Conductance surfaces combine a random
smooth field, habitat mismatch, edge length, water and coast penalties, and a
randomly placed sinuous ridge. Coefficients, length scales, and final edge-log
standard deviation vary among examples. A separate smooth field generates
local $q$. Natural Earth and latent habitat variables generate targets and
sampling designs; they are never supplied to \dreems.

Another 12,256 examples use the $8\times12$ triangular lattice from \feems
simulations, with homogeneous, barrier, and anisotropic surfaces and either
full or approximately 20\% node sampling. A further 5,123 examples use land
or synthetic-corridor graphs with wolf-like sampling paths. The remaining
9,456 use masked grids and earth-like graphs with smooth, barrier,
Gaussian-process, or constant surfaces. Across global-landscape and wolf-like
examples, requested sampled-node fractions are 0.03, 0.05, 0.10, 0.25, or
0.50, weighted toward the sparse values. Each sampled node contains one to
eight diploid individuals, up to 384 total in those families. Other graph
families extend the realized range to 1,200 samples.

\begin{table}[t]
\centering
\caption{Realized production training corpus. Counts refer to graphs, not
SNPs.}
\label{tab:training-corpus}
\small
\begin{tabular}{ll}
\toprule
Corpus property & Value \\
\midrule
Total & 61,440 \\
Global landmass graphs & 34,605 \\
FEEMS triangular lattices & 12,256 \\
Wolf-like graphs & 5,123 \\
Masked grids and earth-like graphs & 9,456 \\
Nodes & 9-1,743 (median 450) \\
Edges & 12-4,107 (median 1,255) \\
Observed-node fraction & 0.002-1.000 (median 0.091) \\
Diploid samples & 2-1,200 (median 200) \\
Accepted SNPs & 2,000-10,000 (median 5,994) \\
Edge log-rate SD & 0.350-1.600 (median 0.976) \\
Node log-q SD & 0.000-0.900 (median 0.452) \\
Simulator & Rust structured coalescent \\
Ascertainment & random marker conditional on MAF >=0.05 \\
\bottomrule
\end{tabular}

\end{table}

\paragraph{Structured-coalescent SNP simulation.}
We wrote a Rust simulator for large stepping-stone models. Its state consists
of ancestral lineages and their current demes. If deme $i$ contains $k_i$
lineages, coalescence occurs there at total rate
$q_i k_i(k_i-1)/2$; each lineage moves from $i$ to $j$ at the specified edge
migration rate. Production rates are symmetric. Descendant sets are stored in
one or more 128-bit words, and random draws use Rust \texttt{SmallRng}.

For the production \texttt{conditioned\_random\_marker} scheme, we simulate a
genealogy and select a branch proportional to branch length, which is the
location distribution of a random infinite-sites mutation conditional on a
mutation occurring on that genealogy. We retain a marker only when its sample
allele frequency lies in $[0.05,0.95]$ and repeat until reaching the requested
number. Within demes, derived and ancestral haplotypes are randomly paired
without replacement to form diploid dosages. Accepted dosage columns are
standardized and accumulated directly into
Equation~\ref{eq:sample-covariance}; genotype matrices are not stored for the
training corpus. Examples contain 2,000--10,000 accepted SNPs.

We compared the Rust simulator with \texttt{msprime} on small graphs and with
symmetry and closed-form checks for the symmetric island model
\citep{baum2022efficient}. Tests compared allele counts, pairwise differences,
and covariance summaries over replicates.

\subsection{Optimization and checkpoint selection}

We used Adam with learning rate $5\times10^{-5}$, minibatch size 8, and global
gradient-norm clipping at one \citep{kingma2015adam}. Training sampled graphs
repeatedly rather than making one pass through the corpus. We evaluated all
1,024 validation graphs every 1,000 steps. After a 20,000-step warmup, early
stopping required an improvement greater than $10^{-4}$ within 12 validation
checks. The best graph-balanced validation loss was 0.15724 at step 29,000.
Training stopped at step 41,000 and restored the step-29,000 parameters. The
one-GPU run took 7,578.6 seconds (2.11 h), used 2,230.1 CPU-seconds, and reached
5.35 GiB resident memory. The checkpoint SHA-256 is
\texttt{f0b677d2cc028252002d15fac843ca7347b1b352eb2286a18bd202b724aa92b4}.

\subsection{Fixed q-aware trust-region covariance refinement}

The neural prediction is amortized and need not closely fit one dataset's
covariance. The default pipeline therefore applies a fixed 50-step refinement
without selecting a test-specific hyperparameter. Let
$P_o=I_o-\vect1\vect1^\mathsf{T}/o$ and let $G_o(\vect\theta)$ be the
observed-node submatrix of $L(\vect\theta)^+$. Define
\begin{equation}
    \bar C_o=P_oC_oP_o,
    \qquad
    \bar G_o(\vect\theta)=P_oG_o(\vect\theta)P_o.
\end{equation}
From the predicted node field we also form
\begin{equation}
    Q_o(\vect\eta)
    =\operatorname{diag}
      \left\{\frac{\exp(\eta_i)}{n_i}:i\text{ observed}\right\},
    \qquad
    \bar Q_o=P_oQ_oP_o,
\end{equation}
where $n_i$ is the number of samples at observed node $i$.

At every edge update, nonnegative least squares fits global scales
$\alpha,\beta\geq0$ for the migration and local-$q$ bases. The objective is
\begin{equation}
 \begin{split}
    \mathcal J(\vect\theta)
    ={}&
    \frac{1}{2\binom{o}{2}}
    \sum_{i<j}
    \left\{
      \bar C_{o,ij}
      -\alpha\bar G_{o,ij}(\vect\theta)
      -\beta\bar Q_{o,ij}
    \right\}^2\\
    &+\frac{10^{-4}}{2m}
      \vect\theta^\mathsf{T}L_E\vect\theta.
 \end{split}
 \label{eq:refinement-objective}
\end{equation}
Analytic gradients through sparse grounded solves avoid constructing $L^+$.
We use 50 Adam updates at learning rate 0.15. Edge batches of 256 control
working memory only.

The update remains in an uncertainty-scaled trust region around the neural
mean $\vect\mu$:
\begin{equation}
    \left[
      \frac{1}{m}\sum_{e=1}^m
      \left\{
        \frac{\theta_e-\mu_e}{\max(\sigma^{\mathrm{raw}}_e,0.25)}
      \right\}^2
    \right]^{1/2}
    \leq0.5.
    \label{eq:trust-region}
\end{equation}
The radius uses the raw neural SD because it was fixed before post-hoc
calibration. The locked configuration has no Gaussian-prior penalty; its
weight is zero. We therefore call the procedure constrained covariance
refinement rather than MAP. Refinement changes the point estimate but not the
reported predictive SD.

\subsection{Synthetic head-to-head benchmark}

\paragraph{Test examples.}
We generated 512 independent global-landscape examples with new random seeds.
Graphs contain 219--1,406 nodes (median 452) and 590--4,065 edges (median
1,270). They contain 2--384 observed nodes (median 39), 5--384 diploid samples
(median 135), and 2,003--9,996 accepted SNPs (median 5,811). The observed-node
fraction ranges from 0.004 to 0.500 (median 0.049). We evaluate both the frozen
neural mean and the fixed radius-0.5 refined mean.

\paragraph{Documented \feems LOO cross-validation.}
We pinned upstream \feems 2.0 at commit
\texttt{aa890b0e8ef583dd657a6a696af591b732d1d8be}. Following its documented
cross-validation notebook and source, each fold holds out all samples at one
observed deme. Candidate fits use the remaining demes, and selection scores
held-out allele-frequency prediction by L2 error. We jointly search ten
migration-smoothness values logarithmically spaced from 100 to 0.001 and five
$q$-smoothness values from 100 to 0.01. The selected pair is refit to all
samples with node-specific $q$, unbounded log weights, and at most 15,000
iterations. This gives \feems its documented leave-one-observed-deme-out,
dataset-specific selector; \dreems uses one checkpoint and one prespecified
refinement for every example. Both methods receive the same genotype matrix,
sample locations, and graph.

\paragraph{Accuracy and paired inference.}
For each graph, we center estimates and true edge log-conductances and compute
Pearson correlation and mean squared error (MSE). Each graph contributes one
value, regardless of edge count. For paired contrasts, an advantage is
\dreems correlation minus \feems correlation or \feems MSE minus \dreems MSE.
We report the mean advantage, a two-sided 95\% Student-$t$ interval, the paired
$t$-test $P$ value, and the fraction of examples won. Sampling strata are
sparse ($\leq10\%$ of graph nodes observed), intermediate ($>10\%$ and
$\leq25\%$), and dense ($>25\%$). Graph-size strata were fixed before
summarization.

For \dreems only, we also evaluate covariance fit by correlating observed
off-diagonal centered covariance with $\alpha\bar G_o+\beta\bar Q_o$ and by
its $R^2$ and MSE. Comparable covariance predictions were not retained in the
completed synthetic \feems artifacts, so these values are diagnostic rather
than a cross-method covariance comparison.

\subsection{Runtime and resource accounting}

Runtime includes model loading, feature construction, the neural forward pass,
and refinement for \dreems, and all LOO folds plus the selected final refit for
\feems. In 511 cases, \dreems used two CPU threads and \feems executed folds
through a process pool of up to 16 workers. The remaining example had 384
\feems folds, each represented as an independent Snakemake job
\citep{koster2012snakemake}; we report its idealized compute critical path and
aggregate worker use separately. These measurements describe the executed
workflows, not hardware-normalized single-core benchmarks. Training and
corpus-generation costs are not charged to an individual \dreems analysis.

\subsection{Predictive-SD calibration and utility}

We used three independent sets of 1,024 simulated graphs for calibration fit,
model selection, and verification. Candidate calibrators left raw SDs
unchanged, multiplied all SDs by one global scale, or applied additional
sampling-density factors for fractions $\leq0.10$, 0.10--0.25, and $>0.25$.
Parameters were fit by graph-balanced normal NLL, and the candidate with
lowest NLL on the selection set was locked. The selected density calibrator
had global multiplier 1.3117 and density factors 1.0753, 0.9951, and 1.9091.
We then measured graph-balanced 50\%, 80\%, and 95\% marginal coverage on the
verification set. The same locked calibration was used in the 512-example
head-to-head benchmark.

To ask whether uncertainty is useful beyond coverage, we correlated calibrated
edge SD with absolute refined error, computed within-graph rank correlations,
correlated graph-mean SD with graph MSE, and compared MSE across pooled SD
quintiles. These are marginal diagnostics for a factorized predictor, not
simultaneous coverage statements.

\subsection{Empirical reanalyses}

We applied the checkpoint without retraining to 15 public datasets that had
georeferenced genotypes, covariance-like EEMS input, or an existing
EEMS/\feems analysis. For raw genotypes we retained biallelic variants,
required study-appropriate call rate, mean-imputed missing calls, and applied a
common minor-allele-frequency threshold of 0.05. Exact source, filtering,
graph construction, and SNP-count provenance is stored per dataset in the
release manifest. Graphs contain 193--4,371 nodes and 539--12,829 edges;
observed fractions range from 0.0066 to 0.2045.

Published comparison panels were copied from the original papers or
author-distributed EEMS output. They can differ from our reanalysis in input
filtering, graph construction, regularization, method, and rendering. We use
them only for a qualitative visual audit. Because true conductance fields are
unknown in empirical data, no edgewise accuracy is reported.

The North American gray-wolf data distributed with \feems received a separate
same-graph analysis \citep{schweizer2016wolves,marcus2021feems}. The common-MAF
filtered matrix contains 111 samples and 17,729 SNPs assigned to 94 of 1,207
nodes on the packaged 3,458-edge graph. We compare neural and refined \dreems
with \feems selected by its 94-fold joint LOO procedure. As descriptive
diagnostics, we report off-diagonal covariance fit, correlation of each
\dreems edge field with the \feems field, edge-field standard deviation, and
mean centered edge log-conductance in a broad Hudson Bay bounding region. The
last quantity is not a formal geographic hypothesis test.
